\documentclass[10pt,a4paper]{amsart}
\usepackage[margin=19mm]{geometry}
\usepackage{amsmath,amssymb,mathtools}
\usepackage[scr=rsfs,cal=euler]{mathalfa}
\usepackage{xfrac}
\usepackage[unicode,hidelinks,bookmarks=false]{hyperref}
\newcommand{\ie}{i.e.\ }
\newcommand{\cf}{cf.\ }
\newcommand{\resp}{respectively\ }
\newcommand{\Subsection}{\subsection}
\providecommand{\nobreakdash}{\nobreak\hbox{-}}
\setlength{\emergencystretch}{2em}
\multlinegap0pt
\usepackage{amssymb}
\usepackage{cancel}
\newtheorem{proposition}{Proposition}
\newtheorem{thm}{Theorem}
\title{Splitting of operations for Hom-diassociative Algebras and Hom-triassociative Algebras}
\author{Abdelkader Hamdouni, Imed Basdouri, Mariem Jendoubi, Ahmed Zahari Abdou Damdji}
\date{Source: arXiv:2601.00040v2 (CC BY 4.0)}
\begin{document}
\maketitle

\noindent\emph{Source and licence.} Splitting of operations for Hom-diassociative Algebras and Hom-triassociative Algebras. Version arXiv:2601.00040v2 (CC BY 4.0), distributed by arXiv under the Creative Commons Attribution 4.0 International licence, http://creativecommons.org/licenses/by/4.0/, as recorded in the arXiv licence metadata for this exact version and shown on the abstract page https://arxiv.org/abs/2601.00040v2. Full licence notice: \texttt{licenses/arxiv-2601.00040-CC-BY-4.0.txt}.\par\smallskip

\noindent\textit{Editorial scope.} Counted unit: the article's thm thm (source0.tex lines 751--761). The article's complete original proof of it is delivered with it (source0.tex lines 762--806, one \texttt{proof} environment of 45 lines). No mathematics is written, repaired or re-derived: the statement and the proof are the article's own lines, in the article's own notation, and no retained line is substituted or re-flowed.  Retained uncounted as the source-local context the proof needs: lines 525--532; lines 590--600; lines 653--671.  Nothing was omitted from any retained span, so the package carries no omission record.  No retained line contains a citation, so the wrapper carries no bibliography and no bibliography entry.  No retained line contains a figure environment, an image-inclusion command or drawing code, so no image is delivered and the package declares no asset dependency.  Editorial formatting only, in addition: an AMS \texttt{amsart} preamble that restates the article's own declarations and macros used by the retained lines, namely the environments \texttt{proposition}, \texttt{thm} on separate counters, together with the standard packages those lines need.  The retained environments print in this document as Proposition 1, Theorem 1 in order of appearance, whereas the article prints the same items with the numbering of its own full document.  The complete native source of the article as distributed in the arXiv e-print for this exact version is bundled unchanged as \texttt{sources/191968/source0.tex}.
\begin{proposition}\label{prop3.11}
Let  $T : \ V \to D$ be a relative averaging operator on a Hom-dendriform algebra \\$(D ,\ \prec,\  \succ ,\ \alpha)$ with respect to the representation $(V ,\  \prec_{\ell} ,\ \succ_{\ell},\  \prec_{r}  ,\ \succ_{r} ,\ \beta)$. Then $V$ inherits a Hom-quadri-dendriform algebra structure with the operations:
\begin{align*}
u \prec^{T}_{\vdash} v &= T(u) \prec_{\ell} v ,\quad u\prec^{T}_{\dashv} v = u \prec_{r} T(v) \\
u\succ^{T}_{\vdash} v &= T(u) \succ_{\ell} v ,\quad u \succ^{T}_{\dashv} v = u \succ_{r} T(v) 
\end{align*}
for all $u ,\ v \in V$.  Moreover  $T$  is  a  homomorphism from  the Hom-quadri-dendriform  algebra \\$(V ,\ \prec^{T}_{\vdash} ,\ \prec^{T}_{\dashv},\  \succ^{T}_{\vdash},\  \succ^{T}_{\dashv},\  \beta)$ to  the  Hom-dendriform algebra $(D ,\  \prec  ,\ \succ ,\ \alpha).$
\end{proposition}

\begin{align}
(x \prec_l v) \prec' \alpha'(w)& = \alpha(x) \prec_l (v \prec' w + v \succ' w) \label{eqtt}\\
(x \succ_l v) \prec' \alpha'(w)& = \alpha(x) \succ_l (v \prec' w) \label{rq2} \\
\alpha(x) \succ_l (v \succ' w)& = (x \prec_l v + x \succ_l v) \succ' \alpha'(w) \label{rq3}\\
(u \prec_r y) \prec' \alpha'(w)& = \alpha'(u) \prec' (y \prec_l w + y \succ_l w) \label{rq4} \\
(u \succ_r y) \prec' \alpha'(w)&= \alpha'(u) \succ' (y \prec_l w) \label{rq5}\\
\alpha'(u)\succ' (y \succ_l w)& = (u \prec_r y + u \succ_r y) \succ' \alpha'(w) \label{rq6}\\
(u \prec' v) \prec_r \alpha(z)& =\alpha'(u) \prec'  (v \prec_r z + v \succ_r z) \label{rq7}\\
(u \succ' v) \prec_r \alpha(z)&= \alpha'(u) \succ' (v \prec_r z) \label{rq8}\\
\alpha'(u) \succ' (v \succ_r z)&= (u \prec' v + u \succ' v) \succ_r \alpha(z).\label{eqttc}
\end{align}

\begin{align}
(x \prec_\vdash y) \prec_\perp \alpha(z) &= \alpha(x) \prec_\vdash (y \prec_\perp z + y \succ_\perp z) \label{sq1}  \\
(x \succ_\vdash y) \prec_\perp \alpha(z) &= \alpha(x) \succ_\vdash (y \prec_\perp z) \label{sq2}\\
\alpha(x) \succ_\vdash (y \succ_\perp z) &= (x \prec_\vdash y + x \succ_\vdash y) \succ_\perp\alpha(z) \label{sq3}\\
(x \prec_\dashv y) \prec_\perp \alpha( z) &= \alpha(x) \prec_\perp (y \prec_\vdash z + y \succ_\vdash z) \label{sq4}\\
(x \succ_\dashv y) \prec_\perp \alpha(z) &= \alpha(x)\succ_\perp (y \prec_\vdash z) \label{sq5}\\
\alpha(x)\succ_\perp (y \succ_\vdash z) &= (x \prec_\dashv y + x \succ_\dashv y) \succ_\perp \alpha(z) \label{sq6}\\
(x \prec_\perp y) \prec_\dashv \alpha(z) &= \alpha(x) \prec_\perp (y \prec_\dashv z + y \succ_\dashv z) \label{sq7} \\
(x \succ_\perp y) \prec_\dashv \alpha(z) &= \alpha(x) \succ_\perp (y \prec_\dashv z) \label{sq8}\\
\alpha(x)\succ_\perp (y \succ_\dashv z) &= (x \prec_\perp y + x \succ_\perp y) \succ_\dashv \alpha(z) \label{sq9}\\
(x \prec_\perp y) \prec_\vdash \alpha(z) &= (x \prec_\vdash y) \prec_\vdash \alpha(z) = (x \prec_\dashv y) \prec_\vdash \alpha(z) \label{sq10}\\
(x \succ_\perp y) \prec_\vdash\alpha(z) &= (x \succ_\vdash y) \prec_\vdash \alpha(z) = (x \succ_\dashv y) \prec_\vdash \alpha(z) \label{sq11}\\
(x \prec_\perp y) \succ_\vdash\alpha(z) &= (x \prec_\vdash y) \succ_\vdash \alpha( z) = (x \prec_\dashv y) \succ_\vdash\alpha(z) \label{sq12}\\
(x \succ_\perp y) \succ_\vdash \alpha(z) &= (x \succ_\vdash y) \succ_\vdash\alpha(z) = (x \succ_\dashv y) \succ_\vdash\alpha(z) \label{sq13}\\
\alpha(x) \prec_\dashv (y \prec_\perp z) &=\alpha(x) \prec_\dashv (y \prec_\vdash z) = \alpha(x) \prec_\dashv (y \prec_\dashv z) \label{sq14}\\
\alpha(x) \succ_\dashv (y \prec_\perp z) &=\alpha(x) \prec_\dashv (y \prec_\vdash z) = \alpha(x) \prec_\dashv (y \prec_\dashv z) \label{sq15}\\
\alpha(x) \succ_\dashv (y \succ_\perp z) &= \alpha(x) \succ_\dashv (y \succ_\vdash z) = \alpha(x) \succ_\dashv (y \succ_\dashv z) \label{sq16}\\
\alpha(x) \prec_\dashv (y \succ_\perp z) &= \alpha(x) \prec_\dashv (y \succ_\vdash z) = \alpha(x) \prec_\dashv (y \succ_\dashv z).\label{sq17}
\end{align}

\begin{thm}
Let $T :\ D' \to D$ be a homomorphic relative averaging operator on the Hom-dendriform algebra $(D  ,\ \prec  ,\ \succ  ,\alpha)$ with respect to an action $(D',\  \prec' ,\ \succ',\  \prec_\ell ,\ \succ_\ell,\  \prec_r ,\ \succ_r ,\ \alpha')$. Then \\$(D',\  \prec_\perp,\  \succ_\perp,\  \prec^T_\vdash,\  \prec^T_\dashv ,\ \succ^T_\vdash ,\ \succ^T_\dashv,\  \alpha')$is a Hom-six-dendriform algebra, where
\begin{align*}
u \prec_\perp v &= u \prec' v ,\quad 
u \succ_\perp v = u \succ' v ,\\ 
u \prec^T_\vdash v &= T u \prec_l v ,\quad
u \prec^T_\dashv v = u \prec_r T v ,\\
u \succ^T_\vdash v &= T u \succ_l v ,\quad 
u \succ^T_\dashv v = u \succ_r T v 
\end{align*}
\end{thm} 

\begin{proof}
We have already $(D',\  \prec_\perp,\  \succ_\perp,\  \alpha')$ is a Hom-dendriform algebra. Moreover  it following from Proposition \ref{prop3.11} that $(D',\  \prec^{T}_\vdash,\   \prec^T_\dashv,\  \succ^T_\vdash,\  \succ^T_\dashv,\  \alpha')$ 
is a Hom-quadri-dendriform algebra. Now  for any $u,\  v,\  w\in D'$ according to Eqs.(\ref{eqtt})-(\ref{eqttc}), we have
\begin{align*}
(u \prec^T_{\vdash} v) \prec_{\perp} \alpha'(w)
&= (Tu \prec_{\ell} v) \prec' \alpha'(w) 
= \alpha (Tu) \prec_{\ell}(v \prec_{\perp} w + v \succ_{\perp} w) \\
&= T\alpha'(u) \prec_{\ell} (v \prec_{\perp} w + v \succ_{\perp} w) 
= \alpha'(u) \prec^T_{\vdash} (v \prec_{\perp} w + v \succ_{\perp} w) \\
(u \succ^T_{\vdash} v) \prec_{\perp} \alpha'(w)
&= (Tu \succ_{\ell} v) \prec' \alpha'(w) 
= \alpha(Tu) \succ_{\ell} (v \prec' w) \\
&= T\alpha'(u) \succ_{\ell} (v \prec' w) 
= \alpha'(u) \succ^T_{\vdash} (v \prec' w) 
= \alpha'(u) \succ^T_{\vdash}(v \succ_{\perp} w) \\
\alpha'(u) \succ^T_{\vdash}(v \succ_{\perp} w)
&= T\alpha'(u) \succ_{\ell} (v \prec' w) 
= (Tu \prec_{\ell} v + Tu \succ_{\ell} v) \succ' \alpha'(w) \\
&= (u \succ^T_{\vdash} v + u \prec^T_{\vdash} v) \succ_{\perp} \alpha'(w) \\
(u \prec^T_{\dashv} v) \prec_{\perp} \alpha'(w) 
&= (u \prec_{r} Tv) \prec' \alpha'(w) 
= \alpha'(u) \prec' (Tv \prec_{\ell} w + Tv \succ_{\ell} w) \\
&= \alpha'(u) \prec_{\perp} (v \prec^T_{\vdash} w + v \succ^T_{\vdash} w) \\
(u \succ^T_{\dashv} v) \prec_{\perp} \alpha'(w) 
&= (u \succ_{r} Tv) \prec' \alpha'(w) 
= \alpha'(u) \succ' (Tv \prec_{\ell} w) 
= \alpha'(u) \succ_{\perp} (v \prec^T_{\vdash} w) \\
\alpha'(u) \succ_{\perp} (v \succ^T_{\vdash} w) 
&= \alpha'(u) \succ' (Tv \succ_{\ell} w) 
= (u \prec_{r} Tv + u \succ_{r} Tv) \succ' \alpha'(w) \\
&= (u \prec^T_{\vdash} v + u \succ^T_{\vdash} v) \succ_{\perp} \alpha'(w) \\
(u \prec_{\perp} v) \prec^T_{\vdash} \alpha'(w) 
&= (u \prec' v) \prec_{r} T\alpha'(w) 
= (u \prec' v) \prec_{r} \alpha(Tw) \\
&= \alpha'(u) \prec' (v \prec_{r} Tw + v \succ_{r} Tw) 
= \alpha'(u) \prec_{\perp} (v \prec^T_{\vdash} w + v \succ^T_{\vdash} w) \\
(u \succ_{\perp} v) \prec^T_{\vdash} \alpha'(w) 
&= (u \succ' v) \prec_{r} T\alpha'(w) 
= (u \succ' v) \prec_{r} \alpha(Tw) \\
&= \alpha'(u) \succ' (v \prec_{r} Tw) 
= \alpha'(u) \succ_{\perp} (v \prec^T_{\vdash} w)
\end{align*}

Similarly it can be proved that other equations  \ref{sq8} à \ref{sq17} are also true. This completes the proof.
\end{proof}

\end{document}
